\documentclass[a4paper, draft, 12pt]{article}
\usepackage{array}
\usepackage{url}
\author{Matthew Yang}
\title{An Extra Boost to Boosted Models}
\date{\today}
\begin{document}
\maketitle
\tableofcontents
\section{Introduction}

Boosted models typically consist of a sequence of decision trees which attempt to predict the previous tree's errors.

Current boosted models are computationally expensive, commonly requiring the utilisation of all the training data at any given time.

Moreover, the structure of current boosted models are largely sequential due to the involvement of decision trees, which significantly slows training speed.

With the ever-growing need to train models on more and more data, we need to look for more computationally efficient alternatives of boosting models.

\section{Aims and Research Questions}

Our research aims to investigate methods of categorising data without an explicit decision tree, whilst still maintaining the structure of a boosted model.

We focus particularly on parallel categorisation techniques, which will be benchmarked against current research efficiency and accuracy.

\section{Methodology}

This method briefly outlines how we will investigate our proposed research.

\subsection{Replication}

We first attempt to replicate the structure of current state-of-the-art models, to build a baseline for our research.

\subsection{Implementation}

We will mimic the structure of current state-of-the-art models, and make appropriate changes where necessary.

\subsection{Benchmarking}

The UCR Time Series Classification Archive~\cite{dau2019ucr} will provide us with a quantitative analysis of our models' performance against other models.

\section{Timeline}

\newcommand{\x}{$\bullet$}

\begin{table}[ht]
\centering
\setlength{\tabcolsep}{5pt}
\renewcommand{\arraystretch}{1.25}
\begin{tabular}{|l|*{12}{c|}}
\hline
Week: & 1 & 2 & 3 & 4 & 5 & 6 & 7 & 8 & 9 & 10 & 11 & 12 \\
\hline
Initial meeting                 & - & \x &    &    &    &    &    &    &    &    &    &    \\
Review of literature            & - &    & \x &    &    &    &    &    &    &    &    &    \\
Writing proposal                & - &    & \x &    &    &    &    &    &    &    &    &    \\
Replication of current research & - &    &    & \x &    &    &    &    &    &    &    &    \\
Implementation of proposal      & - &    &    &    & \x & \x & \x & \x & \x & \x &    &    \\
Benchmarking models             & - &    &    &    &    & \x & \x & \x & \x & \x & \x &    \\
Project presentation            & - &    &    &    &    &    &    &    &    &    &    & \x \\
Final report                    & - &    & \x & \x & \x & \x & \x & \x & \x & \x & \x & \x \\
\hline
\end{tabular}
\caption{Project schedule by week.}
\label{tab:timeline}
\end{table}

\newpage
\bibliographystyle{ieeetr}
\bibliography{references}
\end{document}
